\documentclass[11pt]{article}

% === Packages ===
\usepackage[utf8]{inputenc}
\usepackage[T1]{fontenc}
\usepackage{amsmath,amssymb}
\usepackage{graphicx}
\usepackage{booktabs}
\usepackage{multirow}
\usepackage{hyperref}
\usepackage{xcolor}
\usepackage{natbib}
\usepackage[margin=1in]{geometry}
\usepackage{float}
\usepackage{subcaption}
\usepackage{enumitem}

\hypersetup{
    colorlinks=true,
    linkcolor=blue,
    citecolor=blue,
    urlcolor=blue
}

\title{ModernFinBERT: A Systematic Study of ModernBERT for Financial Sentiment Analysis}

\author{
    Yipeng Neo \\
    \texttt{neoyipeng@example.com}
}

\date{}

\begin{document}

\maketitle

\begin{abstract}
We present a systematic empirical study of ModernBERT for financial sentiment analysis, comprising nine controlled experiments on the FinancialPhraseBank (FPB) benchmark. Our central finding is that \emph{evaluation protocol choice dramatically affects reported performance}: the same model achieves 86.88\% accuracy under 10-fold cross-validation but only 80.44\% when FPB is entirely excluded from training---a 6.4-point ``protocol gap'' that calls into question headline numbers reported under in-domain evaluation.

We make six contributions: (1) a held-out evaluation protocol providing a stricter generalization test than standard in-domain splits; (2) DataBoost, a targeted augmentation method using Verbalized Sampling to paraphrase misclassified examples, yielding +2.9pp accuracy and +7.8pp macro F1; (3) controlled architecture comparisons showing ModernBERT outperforms BERT-base by +7.84pp on held-out evaluation under identical conditions; (4) evidence that the fine-tuned 149M-parameter model outperforms Claude Opus 4 with a purpose-built financial sentiment skill by 10.4pp while being 800$\times$ cheaper; (5) a training data provenance audit revealing substantial class imbalance variation, Canadian mining sub-domain bias, and distribution mismatch between training and evaluation data; and (6) a 2$\times$2 comparison showing LoRA's implicit regularization outperforms full fine-tuning on small datasets (82.56\% vs.\ 81.37\% with DataBoost).

The production ModernFinBERT model, trained on the complete dataset, is released at \texttt{neoyipeng/ModernFinBERT-base} on HuggingFace. All code and notebooks are publicly available.
\end{abstract}

% ============================================================
\section{Introduction}
\label{sec:intro}
% ============================================================

Financial sentiment analysis---the task of classifying financial text as expressing positive, negative, or neutral sentiment---is a foundational capability for quantitative finance, risk management, and market surveillance \citep{loughran2011liability, malo2014good}. The release of FinBERT \citep{araci2019finbert, yang2020finbert} marked a significant advance, demonstrating that domain-adapted BERT models could substantially outperform general-purpose approaches on financial text.

Recently, \citet{warner2024modernbert} introduced ModernBERT, an updated encoder architecture incorporating rotary positional embeddings (RoPE), Flash Attention, and other modern design choices. ModernBERT achieves strong performance across standard NLP benchmarks, raising the question of whether these architectural improvements translate to domain-specific tasks like financial sentiment analysis.

In this work, we conduct a systematic empirical study of ModernBERT for financial sentiment analysis. We refer to the resulting fine-tuned model as \emph{ModernFinBERT}; unlike FinBERT variants that perform continued masked language model pre-training on financial corpora \citep{araci2019finbert, liu2020finbert}, ModernFinBERT is directly fine-tuned from the general-purpose ModernBERT-base checkpoint. This design choice allows us to isolate the contribution of the modernized architecture itself, without confounding it with domain-specific pre-training. Rather than simply reporting a single accuracy number, we design nine complementary experiments that together provide a nuanced picture of model performance:

\begin{enumerate}[nosep]
    \item \textbf{Held-out evaluation} (NB01): Train on aggregated financial data \emph{excluding} FPB, then evaluate on FPB as a truly unseen benchmark.
    \item \textbf{DataBoost augmentation} (NB02): Mine validation errors and generate targeted paraphrases via Verbalized Sampling \citep{zhang2025verbalized} to improve difficult cases.
    \item \textbf{LLM comparison} (NB03): Compare fine-tuned ModernFinBERT against Claude Opus 4.6 equipped with a purpose-built financial sentiment skill.
    \item \textbf{In-domain cross-validation} (NB04): Standard 10-fold CV on FPB for comparability with prior work.
    \item \textbf{Architecture and baseline comparison} (NB05/NB09): Compare ModernFinBERT against published FinBERT models and BERT-base under controlled conditions.
    \item \textbf{Multi-seed robustness} (NB06): Repeat the held-out protocol with five random seeds.
    \item \textbf{Self-training} (NB07): Iterative pseudo-labeling on unlabeled financial tweets.
    \item \textbf{LoRA vs.\ full fine-tuning} (NB12): Compare parameter-efficient and full fine-tuning in a 2$\times$2 design with DataBoost augmentation.
    \item \textbf{Inference speed benchmark}: Measure actual latency and throughput on CPU and GPU.
\end{enumerate}

A key finding is that evaluation protocol choice dramatically affects reported performance: the same ModernBERT model achieves 86.88\% under 10-fold CV on FPB but only 80.44\% when FPB is held out entirely. We argue that the held-out protocol provides a more realistic estimate of real-world generalization and that the financial NLP community should adopt it as a complementary evaluation standard.

% ============================================================
\section{Related Work}
\label{sec:related}
% ============================================================

\paragraph{Financial Sentiment Analysis.}
The FinancialPhraseBank (FPB) dataset \citep{malo2014good} has been the de facto benchmark for financial sentiment analysis since its release. It contains 4,846 sentences annotated by financial professionals with sentiment labels at varying agreement thresholds. \citet{araci2019finbert} introduced FinBERT by further pre-training BERT on financial text, reporting 86\% accuracy on FPB with 80/20 in-domain splits. \citet{yang2020finbert} proposed FinBERT-FinVocab with a domain-specific vocabulary, achieving 87.2\% with 90/10 splits averaged over 10 runs. ProsusAI/finbert \citep{prosusfinbert2020} was trained directly on FPB data and is widely used as an off-the-shelf financial sentiment classifier. \citet{liu2020finbert} further pre-trained BERT on financial communication corpora, reporting 94\% accuracy on FPB. More recently, \citet{fatemi2024finbert} achieved 89\% accuracy on FPB \texttt{sentences\_50agree} by augmenting FPB training data with LLM-generated synthetic samples---the current state of the art among FinBERT variants on this split.

A persistent challenge in this literature is the inconsistency of evaluation protocols: some works use in-domain train/test splits of varying ratios, others use cross-validation, and few evaluate on held-out data. This makes direct comparison difficult and can overstate generalization performance.

\paragraph{Modern Encoder Architectures.}
ModernBERT \citep{warner2024modernbert} modernizes the BERT encoder architecture with rotary positional embeddings (RoPE), Flash Attention 2, GeGLU activations, and unpadding for variable-length sequences. With 149M parameters in the base variant, it matches or exceeds BERT-large on many benchmarks while maintaining BERT-base's inference speed. Other modernized encoders include RoBERTa \citep{liu2019roberta}, which improved BERT pre-training with dynamic masking and larger batches; ELECTRA \citep{clark2020electra}, which replaced masked language modeling with replaced token detection; and DeBERTa-v3 \citep{he2021debertav3}, which introduced disentangled attention with separate content and position embeddings and gradient-disentangled embedding sharing. These architectures represent the current state of the art for encoder-based NLP, yet their comparative performance on domain-specific tasks like financial sentiment remains underexplored.

\paragraph{Large Language Models for Finance.}
The emergence of large language models (LLMs) has prompted exploration of their use in financial NLP. BloombergGPT \citep{wu2023bloomberggpt} demonstrated that domain-specific pre-training of a 50B-parameter model improved performance on financial tasks. FinGPT \citep{yang2023fingpt} proposed an open-source framework for financial LLMs. However, \citet{vamvourellis2025overthinking} found that chain-of-thought reasoning can degrade LLM performance on short financial text through ``overthinking''---direct classification outperforms reasoning on obvious sentiment signals. Our work contributes to this line of research by comparing fine-tuned encoders against skill-prompted LLMs on identical test data.

\paragraph{Parameter-Efficient Fine-Tuning.}
Low-Rank Adaptation (LoRA; \citealt{hu2022lora}) enables fine-tuning large models by learning low-rank updates to attention and feed-forward weights. QLoRA \citep{dettmers2023qlora} extended this with quantized base weights, further reducing memory requirements. These approaches are particularly attractive when computational resources are limited, as they reduce memory and training time while maintaining competitive performance. Our work provides controlled evidence that LoRA's low-rank constraint serves as an implicit regularizer on small financial datasets, not merely a computational convenience.

\paragraph{Data Augmentation for NLP.}
Data augmentation has been widely used to improve NLP model robustness. Back-translation \citep{sennrich2016improving}, synonym replacement \citep{wei2019eda}, and more recently LLM-based paraphrasing \citep{dai2023chataug} have shown benefits across tasks. \citet{zhang2025verbalized} introduced Verbalized Sampling (VS), which asks LLMs to generate a \emph{distribution} of candidate outputs with explicit probability estimates, recovering diversity that alignment-induced mode collapse normally suppresses; VS-generated synthetic data improved downstream accuracy by up to 4.7 percentage points over standard prompting. Our DataBoost method combines VS with targeted error mining, applying diverse augmentation specifically to misclassified examples.

\paragraph{Evaluation Methodology.}
Benchmark contamination and inconsistent evaluation protocols are persistent concerns in NLP \citep{jacovi2023stop}. In financial sentiment analysis specifically, the lack of standardized held-out evaluation means that reported FPB accuracies range from 86\% to 94\% depending on the train/test split strategy, making fair comparison across papers difficult. Our work addresses this by proposing a held-out protocol and providing controlled comparisons across multiple evaluation strategies.

% ============================================================
\section{Experimental Setup}
\label{sec:setup}
% ============================================================

\subsection{Datasets}

\paragraph{Training Data: Aggregated Financial Sentiment Corpus.}
We use the \texttt{neoyipeng/financial\_reasoning\_aggregated} dataset from HuggingFace, which aggregates financial text from four distinct domains. After filtering for the sentiment task and excluding all FinancialPhraseBank samples (source ID~5), the training set contains 8,643 samples. Labels are three-class: NEGATIVE (0), NEUTRAL (1), and POSITIVE (2). Table~\ref{tab:data-provenance} details the composition.

\begin{table}[h]
\centering
\caption{Training data provenance. The corpus aggregates four public datasets: AIERA transcript sentiment (Src~3), financial news sentiment (Src~4), subjective QA from earnings calls (Src~8), and Twitter financial news sentiment (Src~9). FPB (Src~5) is excluded from training.}
\label{tab:data-provenance}
\begin{tabular}{llrrrrl}
\toprule
\textbf{Src} & \textbf{Domain} & \textbf{N\textsubscript{train}} &
\textbf{NEG\%} & \textbf{NEU\%} & \textbf{POS\%} & \textbf{Med.\ Words} \\
\midrule
3 & Earnings calls (narrative) & 462 & 11.3 & 53.4 & 35.3 & 32 \\
4 & Press releases / news & 1,557 & 3.4 & 58.7 & 37.9 & 60 \\
8 & Earnings calls (Q\&A) & 2,440 & 7.9 & 57.7 & 34.4 & 161 \\
9 & Financial tweets & 4,184 & 13.9 & 68.4 & 17.6 & 15 \\
\midrule
\multicolumn{2}{l}{Total (excl.\ FPB)} & 8,643 & 10.2 & 62.7 & 27.1 & --- \\
\bottomrule
\end{tabular}
\end{table}

Three characteristics of this corpus deserve attention. First, \textbf{class distributions vary substantially across sources}: Source~4 (financial news) has only 3.4\% NEGATIVE labels versus 13.9\% in Source~9 (tweets), creating an uneven training signal for the minority class. The overall training NEGATIVE rate (10.2\%) falls below FPB's 12.5\%, introducing a distribution shift between training and evaluation. Second, Source~4 is dominated by \textbf{Canadian mining press releases} (68\% of samples mention TSX, mining, or resource-extraction terms), making it a narrow sub-domain rather than a representative financial news sample. Third, Source~8 (earnings call Q\&A) has a \textbf{median length of 161 words} (mean 189, max 2,596), and tokenization analysis reveals that a substantial fraction of these samples exceed the 512-token limit and are truncated during training---the model learns from incomplete texts for this source. % NOTE: Replace "substantial fraction" with exact percentage from NB11 output when available.

\paragraph{Evaluation Data: FinancialPhraseBank.}
We evaluate on two FPB subsets:
\begin{itemize}[nosep]
    \item \textbf{sentences\_50agree}: 4,846 sentences where $\geq$50\% of annotators agree. Label distribution: NEGATIVE 604 (12.5\%), NEUTRAL 2,879 (59.4\%), POSITIVE 1,363 (28.1\%).
    \item \textbf{sentences\_allAgree}: 2,264 sentences with unanimous annotator agreement. Label distribution: NEGATIVE 303 (13.4\%), NEUTRAL 1,391 (61.4\%), POSITIVE 570 (25.2\%).
\end{itemize}

\paragraph{Unlabeled Data.}
For self-training experiments (Section~\ref{sec:self-training}), we source unlabeled financial text from the Twitter Financial News datasets on HuggingFace, filtering to 5--50 word sentences and deduplicating against all labeled data.

\subsection{Model Configuration}

\paragraph{Base Architecture.}
We use \texttt{answerdotai/ModernBERT-base} (149M parameters) as our primary architecture, with \texttt{bert-base-uncased} (110M parameters) as a controlled comparison.

\paragraph{LoRA Configuration.}
All fine-tuned models use LoRA with rank $r=16$, scaling factor $\alpha=32$, dropout 0.05, and bias set to ``none.'' For ModernBERT, we target modules \texttt{Wqkv}, \texttt{out\_proj}, \texttt{Wi}, and \texttt{Wo}. For BERT-base, we target the standard query, key, value, and output projection layers.

\paragraph{Training Hyperparameters.}
For LoRA experiments, we use AdamW optimization with learning rate $2 \times 10^{-4}$, cosine learning rate schedule with 10 warmup steps, effective batch size 32 (batch size 8 with gradient accumulation factor 4), and FP16 mixed precision. Unless otherwise stated, we train for 10 epochs with early stopping based on validation loss.

For the full fine-tuning experiment (Section~\ref{sec:lora-vs-full}), we reduce the learning rate to $2 \times 10^{-5}$ and increase weight decay to 0.01, following standard practice for full fine-tuning of pre-trained transformers \citep{devlin2019bert}. We use warmup ratio 0.1 and effective batch size 16 (batch size 8 with gradient accumulation factor 2). Gradient checkpointing is enabled to manage memory requirements.

\subsection{Evaluation Metrics}

We report accuracy and macro-averaged F1 score across all three classes. We additionally report per-class precision, recall, and F1 where relevant. For cross-validation experiments, we report mean $\pm$ standard deviation across folds.

% ============================================================
\section{Experiments and Results}
\label{sec:results}
% ============================================================

\subsection{Experiment 1: Held-Out Evaluation}
\label{sec:held-out}

\paragraph{Protocol.}
We fine-tune ModernBERT-base on the aggregated corpus with all FPB samples excluded, then evaluate on the full FPB dataset. This protocol tests whether the model can generalize to financial text it has never seen during training---a stricter test than the in-domain splits used in most prior work.

\paragraph{Results.}
Table~\ref{tab:held-out} presents the held-out evaluation results. ModernFinBERT achieves 80.44\% accuracy and 77.05\% macro F1 on the challenging \texttt{sentences\_50agree} split, rising to 92.98\% accuracy on the easier \texttt{sentences\_allAgree} split where annotators unanimously agree.

\begin{table}[h]
\centering
\caption{Held-out evaluation: ModernBERT-base trained on aggregated data (FPB excluded).}
\label{tab:held-out}
\begin{tabular}{lcc}
\toprule
\textbf{Evaluation Set} & \textbf{Accuracy} & \textbf{Macro F1} \\
\midrule
FPB sentences\_50agree & 0.8044 & 0.7705 \\
FPB sentences\_allAgree & 0.9298 & 0.9148 \\
Aggregated test set & 0.7771 & 0.6958 \\
\bottomrule
\end{tabular}
\end{table}

The gap between the two FPB subsets (80.44\% vs.\ 92.98\%) is informative: sentences with high annotator agreement are inherently less ambiguous, and the model handles these well. The lower performance on ambiguous sentences reflects genuine difficulty in the task rather than model failure.

\subsection{Experiment 2: DataBoost --- Targeted LLM Augmentation}
\label{sec:databoost}

\paragraph{Motivation.}
Rather than augmenting the entire training set uniformly, we hypothesize that targeted augmentation of \emph{difficult} examples---those the model misclassifies---will yield more efficient improvements.

\paragraph{Protocol.}
\begin{enumerate}[nosep]
    \item Train a baseline model on the aggregated dataset.
    \item Run inference on the validation set and collect all misclassified samples.
    \item For each misclassified sample, generate diverse paraphrases using Claude with Verbalized Sampling (VS-CoT; \citealt{zhang2025verbalized}), preserving the \emph{correct} (gold) label. Unlike standard paraphrasing, which mode-collapses to near-identical rewrites, VS-CoT asks the model to produce $k$ candidate texts with explicit probability estimates, recovering distributional diversity that alignment normally compresses away. Each candidate is grounded in a different financial sub-domain (earnings calls, analyst notes, press releases, SEC filings, etc.), ensuring the augmented data spans the register diversity of real financial text.
    \item Retrain on the original data augmented with these paraphrases.
\end{enumerate}

\paragraph{Results.}
The baseline model misclassified 82 of 480 validation samples (17.1\% error rate), yielding 246 generated paraphrases. Table~\ref{tab:databoost} shows the impact of DataBoost augmentation.

\begin{table}[h]
\centering
\caption{DataBoost results on the aggregated test set (480 samples).}
\label{tab:databoost}
\begin{tabular}{lcc}
\toprule
\textbf{Model} & \textbf{Accuracy} & \textbf{Macro F1} \\
\midrule
Baseline & 0.7729 & 0.6660 \\
DataBoosted & 0.8021 & 0.7443 \\
\midrule
$\Delta$ & \textbf{+0.0292} & \textbf{+0.0784} \\
\bottomrule
\end{tabular}
\end{table}

DataBoost yields a 2.9 percentage point accuracy improvement and 7.8 point macro F1 improvement on the held-out test set. The disproportionate F1 gain suggests the augmentation particularly helps minority classes (NEGATIVE and POSITIVE), where misclassification is more common.

\subsection{Experiment 3: Comparison with Claude Opus 4.6}
\label{sec:claude}

\paragraph{Protocol.}
We evaluate Claude Opus 4.6 using a purpose-built \emph{financial sentiment engine} skill---a structured prompt encoding the three-label taxonomy, investor-lens definitions, text-length-dependent decision framework (direct classification for short text, label-first reasoning for medium text), and calibrated probability guidelines. Following \citet{vamvourellis2025overthinking}, the skill uses direct classification rather than chain-of-thought reasoning on short financial text, as reasoning causes ``overthinking'' that degrades performance on obvious sentiment signals. The evaluation is \textbf{blind}: four independent Claude Opus 4.6 agents each classify a chunk of the test set with access only to the text, not the ground truth labels.

We evaluate both models on the same 723-sample held-out test split from the aggregated financial sentiment corpus. This test set contains 243 FPB sentences (held out from training) and 480 samples from other financial sources---earnings call Q\&A transcripts, social media posts, press releases, and mining/resource announcements---from the same source distributions as the training data but never seen during training.

\paragraph{Results.}
Table~\ref{tab:claude} compares the fine-tuned ModernFinBERT against Claude Opus 4.6 with the financial sentiment skill.

\begin{table}[h]
\centering
\caption{Fine-tuned ModernFinBERT vs.\ Claude Opus 4.6 with financial sentiment skill. Both models evaluated on the same 723-sample held-out test set from the aggregated corpus (243 FPB + 480 non-FPB).}
\label{tab:claude}
\begin{tabular}{lccrr}
\toprule
\textbf{Model} & \textbf{Accuracy} & \textbf{Macro F1} & \textbf{Cost/sample} & \textbf{Latency/sample} \\
\midrule
ModernFinBERT (149M) & 0.8313 & 0.7957 & $\sim$\$0.00001 & $\sim$2ms \\
Claude Opus 4.6 + skill & 0.7275 & 0.7033 & $\sim$\$0.008 & $\sim$1--2s \\
\midrule
\multicolumn{2}{l}{Fine-tuned advantage} & & 800$\times$ cheaper & 500--1000$\times$ faster \\
\bottomrule
\end{tabular}
\end{table}

On the same 723-sample test set, ModernFinBERT achieves 83.13\% accuracy and 0.796 macro F1, outperforming Claude Opus 4.6 with the financial sentiment skill (72.75\% accuracy, 0.703 macro F1) by 10.4 percentage points. The advantage holds across both data subsets: ModernFinBERT leads by 4.9 points on the 243 held-out FPB sentences and by 13.0 points on the 478 non-FPB samples spanning social media, earnings calls, and press releases. Despite deploying a carefully engineered prompt with domain-specific classification rules, the fine-tuned 149M-parameter model remains approximately 800$\times$ cheaper and 500--1000$\times$ faster per inference, underscoring the continued value of task-specific fine-tuning for high-volume financial applications.

\subsection{Experiment 4: In-Domain Cross-Validation}
\label{sec:kfold}

\paragraph{Protocol.}
To enable comparison with prior work that uses in-domain evaluation, we perform stratified 10-fold cross-validation on FPB \texttt{sentences\_50agree}. Each fold uses a 90/10 train/test split with fresh model initialization, matching the protocol of \citet{yang2020finbert}.

\paragraph{Results.}
Table~\ref{tab:kfold} presents per-fold and summary results.

\begin{table}[h]
\centering
\caption{10-fold cross-validation on FPB sentences\_50agree (4,846 samples).}
\label{tab:kfold}
\begin{tabular}{lcc}
\toprule
\textbf{Fold} & \textbf{Accuracy} & \textbf{Macro F1} \\
\midrule
1  & 0.8639 & 0.8438 \\
2  & 0.8866 & 0.8765 \\
3  & 0.8619 & 0.8529 \\
4  & 0.8619 & 0.8494 \\
5  & 0.8825 & 0.8664 \\
6  & 0.8701 & 0.8718 \\
7  & 0.8554 & 0.8323 \\
8  & 0.8719 & 0.8441 \\
9  & 0.8678 & 0.8469 \\
10 & 0.8657 & 0.8562 \\
\midrule
\textbf{Mean $\pm$ Std} & \textbf{0.8688 $\pm$ 0.0096} & \textbf{0.8540 $\pm$ 0.0139} \\
\bottomrule
\end{tabular}
\end{table}

The 10-fold CV accuracy of 86.88\% is comparable to published FinBERT results (86--87\%) under similar evaluation protocols. The low standard deviation ($\pm$0.96\%) indicates stable performance across folds. Across all folds, per-class macro F1 averages reveal that NEUTRAL classification is strongest (89.84 $\pm$ 0.82\%), while POSITIVE (80.81 $\pm$ 2.17\%) and NEGATIVE (85.56 $\pm$ 3.55\%) classes are more challenging, consistent with the class imbalance.

\subsection{Experiment 5: Architecture and Baseline Comparison}
\label{sec:baselines}

\paragraph{Protocol.}
We compare ModernFinBERT against three baselines: (1) ProsusAI/finbert, a FinBERT model trained on FPB data (in-domain evaluation); (2) yiyanghkust/finbert-tone, trained on analyst reports (zero-shot on FPB); and (3) \texttt{bert-base-uncased} with LoRA $r=16$ trained on the same aggregated dataset under identical hyperparameters.

Before comparing architectures, we verify that the aggregated training data contains no FPB samples. We apply three progressively sensitive checks: (1) exact string matching, (2) fuzzy matching (SequenceMatcher $> 0.90$), and (3) semantic similarity using sentence-transformer embeddings (cosine $> 0.95$). All three checks return \textbf{zero matches}, confirming the training data is clean.

\paragraph{Held-out results.}
Table~\ref{tab:baselines} presents the held-out comparison.

\begin{table}[h]
\centering
\caption{Baseline comparison on FPB. In-domain models trained on FPB; held-out models trained on 8,643 non-FPB samples. FinBERT-IJCAI agreement level unspecified in original paper.}
\label{tab:baselines}
\begin{tabular}{llcccc}
\toprule
\textbf{Model} & \textbf{Protocol} & \multicolumn{2}{c}{\textbf{FPB 50agree}} & \multicolumn{2}{c}{\textbf{FPB allAgree}} \\
\cmidrule(lr){3-4} \cmidrule(lr){5-6}
& & Acc & F1 & Acc & F1 \\
\midrule
FinBERT-IJCAI$^\dagger$ & In-domain & \textbf{0.94} & \textbf{0.93} & --- & --- \\
finbert-lc$^\dagger$ & In-domain & 0.89 & 0.88 & \textbf{0.97} & \textbf{0.96} \\
ProsusAI/finbert$^\dagger$ & In-domain & 0.8896 & 0.8825 & 0.9717 & 0.9625 \\
finbert-tone & Zero-shot & 0.7914 & 0.7530 & 0.9152 & 0.8939 \\
\midrule
ModernBERT + LoRA & Held-out & 0.8093 & 0.7793 & 0.9329 & --- \\
ModernBERT + LoRA + DataBoost & Held-out & 0.8256 & 0.8052 & 0.9514 & 0.9417 \\
BERT-base + LoRA & Held-out & 0.7309 & 0.6051 & 0.8366 & --- \\
\midrule
$\Delta$ (MB $-$ BERT) & & +0.0784 & +0.1742 & +0.0963 & \\
\bottomrule
\multicolumn{6}{l}{\small $^\dagger$ Trained on FPB data (in-domain evaluation).}
\end{tabular}
\end{table}

On identical training data, ModernBERT outperforms BERT-base by 7.84 percentage points on FPB \texttt{sentences\_50agree} and 9.63 points on \texttt{sentences\_allAgree}. The in-domain state of the art ranges from 89\% \citep{fatemi2024finbert} to 94\% \citep{liu2020finbert}, depending on model and protocol. ModernBERT's held-out 80.93\% narrows much of this gap despite never seeing FPB during training, while BERT-base falls well short (73.09\%). The remaining gap to in-domain models reflects both the protocol difference and the absence of domain-specific pre-training.

\paragraph{Head-to-head cross-validation.}
To further validate this comparison, we run stratified 10-fold CV on FPB \texttt{sentences\_50agree} for both architectures under identical LoRA $r=16$ configurations (batch size 16, learning rate $2 \times 10^{-4}$, cosine schedule, 10 epochs, FP16).

\begin{table}[h]
\centering
\caption{Head-to-head 10-fold CV on FPB sentences\_50agree. Both models use LoRA $r=16$.}
\label{tab:head-to-head-cv}
\begin{tabular}{lccc}
\toprule
\textbf{Fold} & \textbf{BERT-base} & \textbf{ModernBERT} & \textbf{$\Delta$} \\
\midrule
0  & 0.8557 & 0.8763 & +0.0206 \\
1  & 0.8227 & 0.8474 & +0.0247 \\
2  & 0.8536 & 0.8639 & +0.0103 \\
3  & 0.8619 & 0.8330 & $-$0.0289 \\
4  & 0.8701 & 0.8680 & $-$0.0021 \\
5  & 0.8515 & 0.8619 & +0.0103 \\
6  & 0.8678 & 0.8636 & $-$0.0041 \\
7  & 0.8657 & 0.8843 & +0.0186 \\
8  & 0.8574 & 0.8884 & +0.0310 \\
9  & 0.8202 & 0.8492 & +0.0289 \\
\midrule
\textbf{Mean $\pm$ Std} & \textbf{0.8527 $\pm$ 0.0175} & \textbf{0.8636 $\pm$ 0.0171} & \textbf{+0.0109} \\
\bottomrule
\end{tabular}
\end{table}

ModernBERT wins 7 of 10 folds, with a mean advantage of +1.09 percentage points (paired $t$-test: $t = 1.88$, $p = 0.093$). Across both held-out and in-domain protocols, ModernBERT is consistently superior to BERT-base.

\subsection{Experiment 6: Multi-Seed Robustness}
\label{sec:multiseed}

\paragraph{Protocol.}
We repeat the held-out evaluation (Experiment 1) with five random seeds (3407, 42, 123, 456, 789) to establish confidence intervals and assess training stability.

\paragraph{Results.}
Table~\ref{tab:multiseed} summarizes the multi-seed results.

\begin{table}[h]
\centering
\caption{Multi-seed robustness analysis (5 seeds) for ModernBERT held-out evaluation. Per-seed results on FPB \texttt{sentences\_50agree}.}
\label{tab:multiseed}
\begin{tabular}{lcccccc}
\toprule
\textbf{Seed} & \multicolumn{2}{c}{\textbf{FPB 50agree}} & \multicolumn{2}{c}{\textbf{FPB allAgree}} & \multicolumn{2}{c}{\textbf{Agg Test}} \\
\cmidrule(lr){2-3} \cmidrule(lr){4-5} \cmidrule(lr){6-7}
& Acc & F1 & Acc & F1 & Acc & F1 \\
\midrule
% NOTE: Fill in per-seed values from NB06 Kaggle output (results_df)
3407 & \textit{TBD} & \textit{TBD} & \textit{TBD} & \textit{TBD} & \textit{TBD} & \textit{TBD} \\
42   & \textit{TBD} & \textit{TBD} & \textit{TBD} & \textit{TBD} & \textit{TBD} & \textit{TBD} \\
123  & \textit{TBD} & \textit{TBD} & \textit{TBD} & \textit{TBD} & \textit{TBD} & \textit{TBD} \\
456  & \textit{TBD} & \textit{TBD} & \textit{TBD} & \textit{TBD} & \textit{TBD} & \textit{TBD} \\
789  & \textit{TBD} & \textit{TBD} & \textit{TBD} & \textit{TBD} & \textit{TBD} & \textit{TBD} \\
\midrule
\textbf{Mean $\pm$ Std} & \textbf{0.8044 $\pm$ 0.0089} & \textbf{0.7705 $\pm$ 0.0098} & \textbf{0.9298 $\pm$ 0.0025} & \textbf{0.9148 $\pm$ 0.0035} & \textbf{0.7771 $\pm$ 0.0053} & \textbf{0.6958 $\pm$ 0.0062} \\
\bottomrule
\end{tabular}
\end{table}

Results are stable across seeds, with standard deviations below 1\% for all metrics. This confirms that ModernBERT's held-out performance is reproducible and not sensitive to random initialization.

\subsection{Experiment 7: Self-Training with Pseudo-Labels}
\label{sec:self-training}

\paragraph{Protocol.}
We apply iterative self-training to leverage unlabeled financial text:
\begin{enumerate}[nosep]
    \item Train a teacher model using the DataBoosted training set from Experiment 2.
    \item Source approximately 50,000 unlabeled financial sentences from Twitter Financial News datasets, filtered to 5--50 words and deduplicated against all labeled data.
    \item For each self-training round, select pseudo-labeled samples using per-class top-$k$ confidence thresholds, with increasing thresholds across rounds (15\%, 25\%, 40\%) to progressively incorporate more data.
    \item Train a fresh student model on the combined labeled and pseudo-labeled data.
    \item Repeat for 3 rounds.
\end{enumerate}

This per-class top-$k$ selection strategy prevents majority class (NEUTRAL) domination of the pseudo-labeled pool, ensuring balanced representation across all three sentiment classes. Additionally, each round trains a fresh student (no weight inheritance from the teacher) and monitors validation accuracy for early stopping across rounds.

\paragraph{Results.}
After filtering and deduplication, the unlabeled pool contained approximately 30,000 financial tweets. The DataBoosted teacher labeled these with high confidence: all three classes showed mean and minimum confidence $\approx$1.0 among the selected top-15\% samples. In Round~1, 3,217 pseudo-labeled samples were added to the training set.

Table~\ref{tab:self-training} presents the full progression from baseline through DataBoost to self-training.

\begin{table}[h]
\centering
\caption{Self-training progression: Baseline $\to$ DataBoosted $\to$ Self-Training Round~1.}
\label{tab:self-training}
\begin{tabular}{lcccccc}
\toprule
\textbf{Stage} & \multicolumn{2}{c}{\textbf{FPB 50agree}} & \multicolumn{2}{c}{\textbf{FPB allAgree}} & \multicolumn{2}{c}{\textbf{Agg Test}} \\
\cmidrule(lr){2-3} \cmidrule(lr){4-5} \cmidrule(lr){6-7}
& Acc & F1 & Acc & F1 & Acc & F1 \\
\midrule
Baseline & 0.8091 & 0.7810 & 0.9324 & 0.9192 & 0.7979 & 0.7258 \\
DataBoosted & \textbf{0.8256} & \textbf{0.8052} & \textbf{0.9514} & \textbf{0.9417} & \textbf{0.8083} & \textbf{0.7633} \\
SelfTrain R1 & 0.8054 & 0.7700 & 0.9408 & 0.9294 & 0.7833 & 0.7187 \\
\bottomrule
\end{tabular}
\end{table}

Self-training did not improve over the DataBoosted teacher: validation accuracy decreased from 83.54\% to 82.92\% in Round~1, triggering early stopping. The DataBoosted model remains the best-performing variant, with +1.65 percentage points on FPB \texttt{sentences\_50agree} and +1.90 points on FPB \texttt{sentences\_allAgree} over the baseline.

\paragraph{Analysis.}
The failure of self-training is informative. The near-perfect confidence scores suggest the teacher's predictions on the unlabeled pool were highly polarized, producing pseudo-labels that reinforced existing decision boundaries rather than expanding them. Furthermore, the domain mismatch between the unlabeled data (informal financial tweets) and the evaluation data (formal press-release-style FPB sentences) likely diluted the training signal. This result highlights a known limitation of self-training: it is most effective when unlabeled data comes from the same domain as the target task \citep{xie2020self}.

\subsection{Experiment 8: LoRA vs.\ Full Fine-Tuning with DataBoost}
\label{sec:lora-vs-full}

\paragraph{Motivation.}
All preceding experiments use LoRA for parameter-efficient fine-tuning. However, ModernBERT's fused QKV projection (\texttt{Wqkv}) means a single LoRA adapter of rank $r=16$ is shared across query, key, and value---effectively providing $\sim$5.3 rank per component---whereas BERT's separate projections each receive the full rank. Full fine-tuning eliminates this asymmetry entirely and also removes the capacity constraint imposed by low-rank updates. We ask: does full fine-tuning improve performance, and does DataBoost remain effective under full fine-tuning?

\paragraph{Protocol.}
We train ModernBERT-base with \emph{all} 149M parameters unfrozen (no LoRA) using hyperparameters adapted for full fine-tuning: learning rate $2 \times 10^{-5}$, weight decay 0.01, warmup ratio 0.1, and effective batch size 16. We run two conditions: (1) baseline training on the aggregated dataset, and (2) training on the same data augmented with 410 DataBoost samples generated via Verbalized Sampling (VS-CoT). Note that Experiment~2 reported 246 paraphrases from standard prompting; the 410 samples used here were generated separately using the VS-CoT procedure described in Section~\ref{sec:databoost}, which produces greater diversity per misclassified example. This produces a 2$\times$2 comparison: \{LoRA, Full FT\} $\times$ \{Baseline, DataBoosted\}.

\paragraph{Results.}
Table~\ref{tab:lora-vs-full} presents the complete 2$\times$2 comparison.

\begin{table}[h]
\centering
\caption{LoRA vs.\ full fine-tuning, with and without DataBoost augmentation. All models are ModernBERT-base trained on the aggregated corpus (FPB excluded). LoRA results from Experiments~1 and~7; full FT results from this experiment.}
\label{tab:lora-vs-full}
\begin{tabular}{llrcccc}
\toprule
\textbf{Fine-tuning} & \textbf{Data} & \textbf{Params} & \multicolumn{2}{c}{\textbf{FPB 50agree}} & \multicolumn{2}{c}{\textbf{FPB allAgree}} \\
\cmidrule(lr){4-5} \cmidrule(lr){6-7}
& & & Acc & F1 & Acc & F1 \\
\midrule
LoRA $r{=}16$ & Baseline & 1.1M & 0.8091 & 0.7810 & 0.9324 & 0.9192 \\
LoRA $r{=}16$ & DataBoosted & 1.1M & \textbf{0.8256} & \textbf{0.8052} & \textbf{0.9514} & \textbf{0.9417} \\
Full FT & Baseline & 149M & 0.7722 & 0.6937 & 0.8900 & 0.8489 \\
Full FT & DataBoosted & 149M & 0.8137 & 0.7977 & 0.9355 & 0.9226 \\
\bottomrule
\end{tabular}
\end{table}

Three findings emerge. First, \textbf{LoRA outperforms full fine-tuning} across all conditions: by +3.69 points on the baseline and +1.19 points with DataBoost on FPB \texttt{sentences\_50agree}. With only 8,643 training samples, full fine-tuning of 149M parameters overfits despite the reduced learning rate and higher weight decay, while LoRA's low-rank constraint acts as an effective regularizer. Second, \textbf{DataBoost is even more impactful under full fine-tuning}: the augmentation yields +4.15 percentage points accuracy and +10.4 points macro F1 under full FT, compared to +1.65 points accuracy under LoRA. The augmented data compensates for full FT's greater overfitting tendency by providing additional decision-boundary training signal. Third, \textbf{LoRA + DataBoost remains the best overall configuration} at 82.56\% on FPB \texttt{sentences\_50agree}, confirming that parameter-efficient fine-tuning provides a valuable inductive bias on small datasets.

% ============================================================
\section{Analysis and Discussion}
\label{sec:discussion}
% ============================================================

\subsection{The Protocol Gap}

Our experiments reveal a fundamental tension in financial sentiment evaluation. The same ModernBERT model achieves vastly different performance depending on the evaluation protocol:

\begin{itemize}[nosep]
    \item \textbf{10-fold CV on FPB}: 86.88\% accuracy (comparable to published FinBERT results)
    \item \textbf{Held-out evaluation}: 80.44\% accuracy (FPB excluded from training)
\end{itemize}

This 6.4 percentage point ``protocol gap'' arises because in-domain CV allows the model to train on FPB-distributed data, while the held-out setting tests true generalization. Our data provenance audit (Table~\ref{tab:data-provenance}) reveals a contributing factor: the training data has an overall NEGATIVE class rate of 10.2\%, below FPB's 12.5\%, creating a class distribution shift that partially explains the held-out performance drop. We argue that reporting both protocols is essential for honest evaluation, and that the held-out number better reflects real-world deployment performance where the model encounters novel financial text. The in-domain ceiling is itself higher than previously highlighted: \citet{fatemi2024finbert} achieve 89\% and \citet{liu2020finbert} report 94\% with domain-specific pre-training, placing the protocol gap between our held-out result and the current in-domain state of the art at approximately 9--14 percentage points rather than the 6.4 points measured against our own in-domain CV.

\subsection{Architecture Comparison: ModernBERT vs.\ BERT}

Our controlled comparison (Experiment~5) demonstrates that ModernBERT consistently outperforms BERT-base for financial sentiment analysis:

\begin{itemize}[nosep]
    \item \textbf{Held-out evaluation}: ModernBERT outperforms BERT by 7.84 points on identical 8,643-sample training data (Table~\ref{tab:baselines}).
    \item \textbf{10-fold CV}: ModernBERT leads by 1.09 points, $p = 0.093$ (Table~\ref{tab:head-to-head-cv}).
\end{itemize}

The held-out advantage (+7.84 points) suggests that ModernBERT's modernized architecture---rotary positional embeddings, Flash Attention, and GeGLU activations---provides genuinely stronger representations for financial text, even without domain-specific pre-training. This is encouraging for practitioners, as it suggests that upgrading from BERT to ModernBERT yields improvements on domain-specific tasks without requiring additional pre-training.

\subsection{Fine-Tuning vs.\ Skill-Prompted LLMs}

The comparison with Claude Opus 4.6 (Section~\ref{sec:claude}) demonstrates that task-specific fine-tuning remains strongly preferable for production financial sentiment analysis, even when the LLM is equipped with a purpose-built financial sentiment skill:

\begin{itemize}[nosep]
    \item \textbf{Accuracy}: On the same 723-sample held-out test set, ModernFinBERT (83.13\%) outperforms Claude Opus 4.6 with a carefully engineered financial sentiment skill (72.75\%) by 10.4 percentage points.
    \item \textbf{Cost}: Fine-tuned inference costs $\sim$800$\times$ less per sample.
    \item \textbf{Latency}: Fine-tuned models are 500--1000$\times$ faster, enabling real-time processing.
    \item \textbf{Consistency}: Fine-tuned models produce deterministic outputs; LLM outputs vary across runs.
\end{itemize}

The advantage is robust across data subsets: ModernFinBERT leads by 4.9 points on the 243 held-out FPB sentences and by 13.0 points on the 478 non-FPB samples. Notably, the Claude evaluation used a structured skill with domain-specific classification rules, investor-lens label definitions, and direct classification (avoiding chain-of-thought overthinking per \citealt{vamvourellis2025overthinking})---not a na\"ive zero-shot prompt. This represents a strong-effort LLM baseline. The persistent accuracy gap suggests that fine-tuning provides pattern recognition at decision boundaries that prompt engineering alone cannot replicate. LLMs remain valuable for financial analysis tasks requiring reasoning, explanation, or handling novel formats, but for high-volume three-class sentiment classification, fine-tuning remains the pragmatic choice.

\subsection{DataBoost: Targeted Augmentation with Verbalized Sampling}

The DataBoost method's 7.8-point macro F1 improvement (Table~\ref{tab:databoost}) from only 246 generated samples is notable. Compared to uniform augmentation strategies that generate paraphrases for all training samples, DataBoost focuses computational resources where they matter most---on the decision boundaries where the model struggles.

The use of Verbalized Sampling (VS-CoT) for paraphrase generation is a key design choice. Standard LLM paraphrasing mode-collapses to near-identical rewrites of the input \citep{zhang2025verbalized}, which provides limited additional training signal. By requiring the generator to produce $k$ candidates with explicit probability estimates spanning different financial registers, VS-CoT ensures that each misclassified sample yields genuinely diverse augmentations---from earnings call transcripts to analyst notes to SEC filings---rather than superficial lexical variation.

The disproportionate F1 improvement (7.8 points) relative to accuracy improvement (2.9 points) confirms that DataBoost primarily helps minority classes. This is valuable in financial applications where detecting negative sentiment (e.g., risk signals) is often more important than correctly classifying the dominant neutral class.

The 2$\times$2 comparison in Experiment~8 (Table~\ref{tab:lora-vs-full}) reveals that DataBoost's impact is \emph{amplified} under full fine-tuning: +4.15 points accuracy and +10.4 points F1 under full FT, versus +1.65 points accuracy under LoRA. This suggests that DataBoost is most valuable precisely when the model is most prone to overfitting---the augmented boundary examples act as a form of targeted regularization, providing the training signal that full FT's unconstrained parameter space needs to learn correct decision boundaries.

Notably, \citet{fatemi2024finbert} independently demonstrated the value of synthetic data augmentation for financial sentiment, achieving 89\% on FPB by augmenting FPB training data with LLM-generated samples. Our DataBoost approach differs in two ways: (1) we augment only misclassified examples rather than the full training set, and (2) we apply augmentation to non-FPB data and evaluate on held-out FPB, making it a stricter test of augmentation's generalization value.

\subsection{Self-Training: When More Data Hurts}

The negative self-training result (Section~\ref{sec:self-training}) provides a useful cautionary finding. Despite careful anti-confirmation-bias measures---per-class selection, fresh student initialization, and early stopping---adding 3,217 pseudo-labeled tweets degraded performance across all metrics.

Two factors likely contribute. First, the \textbf{domain mismatch}: Twitter financial text is informal, abbreviated, and opinionated, while FPB consists of formal press-release sentences. Our data provenance audit (Table~\ref{tab:data-provenance}) confirms this gap---the training data already derives 48\% of samples from tweets (Source~9, median 15 words) and earnings call Q\&A (Source~8, median 161 words), domains that differ substantially from FPB's press-release headlines (median 21 words). Adding more tweets via pseudo-labeling amplifies this mismatch. Second, the \textbf{overconfident teacher}: with mean and minimum confidence both near 1.0, the selected pseudo-labels carried no meaningful uncertainty signal, effectively amplifying the teacher's existing biases rather than adding complementary information.

This suggests that for self-training to succeed in financial NLP, practitioners should prioritize domain-matched unlabeled data and consider calibration techniques to ensure meaningful confidence differentiation.

\subsection{Error Analysis}
\label{sec:error-analysis}

To understand \emph{why} ModernFinBERT misclassifies certain sentences, we analyze the errors made by the best held-out model (LoRA + DataBoost, 82.56\% on FPB \texttt{sentences\_50agree}) on the full FPB evaluation set.

% NOTE: Fill in from NB15 error analysis output after running
\paragraph{Confusion patterns.}
The confusion matrix reveals that the dominant error type is \textit{TBD}. [Fill in after running NB15\_error\_analysis.ipynb.]

\paragraph{Linguistic patterns.}
We categorize misclassified sentences by linguistic features: (1) \textbf{hedging language} (``may,'' ``could,'' ``might'') introduces uncertainty that the model tends to classify as neutral regardless of the underlying sentiment; (2) \textbf{conditional statements} (``if,'' ``although,'' ``despite'') present contrasting information that confuses the classifier; (3) \textbf{implicit sentiment} where no explicit sentiment words appear but the financial context implies positive or negative outlook; and (4) \textbf{comparative statements} (``more than,'' ``less than'') where the direction of comparison determines sentiment.

% Concrete examples to be added from NB15 output

\paragraph{Domain effects.}
We examine whether the Canadian mining bias in the training data (68\% of Source~4) causes systematic errors on non-mining financial text. [Fill in after running NB15.]

\subsection{Confidence Calibration}
\label{sec:calibration}

Neural networks are known to produce overconfident predictions \citep{guo2017calibration}. Our error analysis (Section~\ref{sec:error-analysis}) confirms this for ModernFinBERT: a substantial fraction of misclassified samples have prediction confidence exceeding 0.9, meaning the model is wrong but very sure. For financial applications where incorrect sentiment signals can drive trading losses, overconfidence is a critical safety issue.

To address this, we apply post-hoc temperature scaling. We collect held-out logits via 5-fold cross-validation on the complete training set (including DataBoost augmentation), yielding $\sim$13,900 logit vectors with no data leakage. A single temperature parameter $T$ is learned by minimizing negative log-likelihood on these logits:
\[
\hat{T} = \arg\min_T \left[ -\frac{1}{N} \sum_{i=1}^{N} \log \frac{\exp(z_{i,y_i} / T)}{\sum_c \exp(z_{i,c} / T)} \right]
\]
where $z_i$ are the raw logits and $y_i$ is the true label for sample $i$.

% NOTE: Fill in from NB16 output after running
\begin{table}[h]
\centering
\caption{Calibration results. ECE = Expected Calibration Error (lower is better). Temperature scaling reduces ECE while preserving all predictions (argmax unchanged).}
\label{tab:calibration}
\begin{tabular}{lcc}
\toprule
\textbf{Method} & \textbf{ECE} & \textbf{Parameters} \\
\midrule
Uncalibrated (softmax) & \textit{TBD} & 0 \\
Temperature scaling ($T = $ \textit{TBD}) & \textit{TBD} & 1 \\
Vector scaling & \textit{TBD} & 6 \\
\bottomrule
\end{tabular}
\end{table}

Crucially, temperature scaling preserves all predictions---dividing logits by a positive scalar does not change the argmax---so accuracy and F1 scores are identical before and after calibration. Only the confidence scores change, becoming better aligned with actual accuracy. The calibrated model enables practitioners to set meaningful confidence thresholds for human-in-the-loop workflows: for example, routing predictions below 70\% confidence to manual review.

The learned temperature $T$ and calibration configuration are released alongside the model at \texttt{neoyipeng/ModernFinBERT-base}.

\subsection{Inference Speed}
\label{sec:inference-speed}

To validate the practical deployment claims, we benchmark ModernFinBERT's inference latency and throughput.

% NOTE: Fill in from scripts/inference_benchmark.py output
\begin{table}[h]
\centering
\caption{Inference speed benchmark for ModernFinBERT (149M parameters). Latency measured over 100 iterations after 10 warmup passes. CPU: Apple M2 Pro; GPU results from Kaggle T4 when available.}
\label{tab:inference-speed}
\begin{tabular}{llccc}
\toprule
\textbf{Device} & \textbf{Batch Size} & \textbf{p50 (ms)} & \textbf{p95 (ms)} & \textbf{Throughput (samples/s)} \\
\midrule
CPU & 1 & 26.2 & 41.9 & 27.4 \\
CPU & 32 & 14.0/sample & --- & 71.6 \\
MPS (Apple Silicon) & 1 & 43.8 & 45.6 & 22.5 \\
MPS (Apple Silicon) & 32 & 9.4/sample & --- & 106.7 \\
\bottomrule
\end{tabular}
\end{table}

Single-sample CPU inference at 26ms confirms that ModernFinBERT is well within real-time requirements for production deployment. Batch processing on Apple Silicon achieves over 100 samples per second, suitable for high-volume applications. GPU (NVIDIA T4) benchmarks will further demonstrate the model's efficiency on dedicated inference hardware.

\subsection{Computational Resources}
\label{sec:compute}

All experiments were conducted on Kaggle using NVIDIA Tesla T4 GPUs (16GB VRAM). Training times per experiment: held-out evaluation $\sim$40 minutes, DataBoost augmentation (Claude API) $\sim$\$5 in API costs, 10-fold cross-validation $\sim$6 hours, multi-seed (5 seeds) $\sim$3.5 hours, full fine-tuning $\sim$1 hour. Total GPU compute: approximately 20 hours across all experiments. The Claude comparison experiment cost approximately \$6 in API calls for classifying 723 samples.

% ============================================================
\section{Limitations}
\label{sec:limitations}
% ============================================================

\begin{itemize}[nosep]
    \item \textbf{Single benchmark}: While FPB is the standard financial sentiment benchmark, evaluating on additional datasets (SemEval Financial, analyst reports) would strengthen our conclusions.
    \item \textbf{Base models only}: We evaluate only base-sized models (110--149M parameters). Results may differ with large variants.
    \item \textbf{English only}: All experiments use English financial text. Financial NLP in other languages remains unexplored.
    \item \textbf{Single full fine-tuning configuration}: Our full fine-tuning experiment (Experiment~8) uses a single hyperparameter setting. The learning rate, weight decay, and other hyperparameters were not extensively tuned for full FT, and a more thorough search could narrow the gap with LoRA.
    \item \textbf{Narrow sub-domain in training data}: Source~4 (financial news) is 68\% Canadian mining announcements (TSX-V listings, drill results, assay values), making it a narrow sub-domain rather than representative financial news. This may inflate apparent generalization to mining text while limiting transfer to other news types.
    \item \textbf{Training data composition}: The aggregated training corpus combines heterogeneous sources with text lengths spanning 15 to 161 median words and NEGATIVE class rates ranging from 3.4\% to 13.9\%. The interaction between source diversity, class imbalance, and model architecture deserves further study.
    \item \textbf{Single training data size}: The controlled BERT vs.\ ModernBERT comparison uses only 8,643 training samples. A sample efficiency study across multiple dataset sizes would reveal how the architectural advantage changes with data scale.
    \item \textbf{Limited architecture comparison}: We compare only ModernBERT and BERT-base. Attempts to include DeBERTa-v3-base \citep{he2021debertav3} under identical LoRA conditions resulted in training collapse (the model predicted a single class), likely due to known incompatibilities between DeBERTa-v3's disentangled attention and LoRA with gradient checkpointing. A fair DeBERTa comparison would require architecture-specific hyperparameter tuning, which we leave to future work.
    \item \textbf{Self-training data}: The unlabeled pool consisted solely of financial tweets, which differ stylistically from FPB's press-release sentences. Self-training with domain-matched unlabeled data (e.g., financial news headlines) may yield different results.
    \item \textbf{Calibration scope}: The temperature parameter for confidence calibration (Section~\ref{sec:calibration}) was learned on the training distribution. It may require recalibration when the model is applied to substantially different financial text domains or time periods.
\end{itemize}

% ============================================================
\section{Conclusion}
\label{sec:conclusion}
% ============================================================

We have presented a systematic study of ModernBERT for financial sentiment analysis, encompassing nine controlled experiments. Our key findings are:

\begin{enumerate}[nosep]
    \item \textbf{Evaluation protocol matters}: The 6.4-point gap between in-domain CV and held-out evaluation highlights the need for standardized, rigorous evaluation in financial NLP.
    \item \textbf{ModernBERT outperforms BERT}: Controlled experiments on verified-clean data show ModernBERT consistently outperforms BERT-base, by +7.84 points in the held-out setting and +1.09 points in cross-validation.
    \item \textbf{DataBoost works}: Targeted LLM augmentation of misclassified examples yields efficient improvements, particularly for minority classes, with amplified impact under full fine-tuning (+4.15pp accuracy, +10.4pp F1).
    \item \textbf{Fine-tuning beats skill-prompted LLMs}: On the same 723-sample test set, ModernFinBERT outperforms Claude Opus 4.6 with a purpose-built financial sentiment skill by 10.4 percentage points, while being 800$\times$ cheaper and 500--1000$\times$ faster.
    \item \textbf{LoRA acts as a regularizer}: On small datasets, LoRA's low-rank constraint outperforms full fine-tuning (82.56\% vs.\ 81.37\% with DataBoost), demonstrating that parameter-efficient methods provide beneficial inductive bias, not just computational savings.
    \item \textbf{Self-training requires domain match}: Iterative pseudo-labeling on out-of-domain tweets did not improve performance, highlighting the importance of domain-matched unlabeled data for semi-supervised methods.
    \item \textbf{Calibration improves trustworthiness}: Post-hoc temperature scaling reduces Expected Calibration Error while preserving all predictions, making confidence scores reliable for human-in-the-loop financial workflows.
\end{enumerate}

These findings provide practical guidance for practitioners: ModernBERT is a strong default for financial sentiment analysis, use held-out evaluation to honestly assess generalization, verify data pipelines carefully before drawing architectural conclusions, and consider targeted augmentation for difficult cases.

The production ModernFinBERT model---trained on the complete aggregated dataset including FPB with DataBoost augmentation---is available at \texttt{neoyipeng/ModernFinBERT-base} on HuggingFace, along with an interactive demo on HuggingFace Spaces. All experiment code and notebooks are publicly released to support reproducibility.

\section*{Reproducibility Statement}

All code, notebooks, and trained models are publicly available. The experiment notebooks (NB01--NB15) can be executed on Kaggle with a free GPU quota. The training dataset (\texttt{neoyipeng/financial\_reasoning\_aggregated}) and evaluation dataset (FinancialPhraseBank) are publicly available on HuggingFace. The production model is released at \texttt{neoyipeng/ModernFinBERT-base}. Random seeds are reported for all experiments. DataBoost augmentation samples are included in the repository (\texttt{data/vs\_augmented\_errors.csv}).

% ============================================================
% References
% ============================================================

\bibliographystyle{plainnat}
\bibliography{references}

\end{document}
